\begin{auxlemma}[Orbit transference from the finite estimate]\label{l:orbit-transference}
\uses{t:teravainen,def:averages}
\lean{PSPrimes.OrbitVNInput, PSPrimes.orbitVNInput_of_finiteVNInput}\leanok
Assume TV. For each maximum degree $d\geq1$, the same degree-only exponent
$K\geq1$ works as follows. Fix integer polynomials $P_i$ of distinct positive
increasing degrees, largest degree $d$. There is $C_{\mathbf P}>0$, chosen
before the system, its functions and $N$, such that every invertible probability
system, essentially $1$-bounded a.e. strongly measurable $f_i$, $N\geq1$, and
real weight $g$ satisfying $|g(n)|\leq B$, $B>0$, and
\[
 \left|\sum_{n=1}^N g(n)\ee(\phi(n))\right|\leq NE
 \quad(\deg\phi\leq d),\qquad E\geq0,
\]
satisfy integrability and the bound
\[
 \int_X\left|\frac1N\sum_{n=1}^N g(n)
     \prod_i f_i(T^{P_i(n-1)}x)\right|\,d\nu(x)
 \leq C_{\mathbf P}B\bigl(N^{-1}+E/B\bigr)^{1/K}.
\]
\end{auxlemma}
\Status{checked derivation from TV of the required real-weight specialization of the related manuscript's Proposition 1.4; no independent own-paper premise}
\begin{proof}[Duality, exact spatial averaging and finite cutoffs]
\proves{l:orbit-transference}\uses{t:teravainen,def:averages}\leanok
For exposition first suppose $|f_i|\leq1$ everywhere and put
\[
 A(x)=N^{-1}\sum_{n=1}^N g(n)\prod_i f_i(T^{P_i(n-1)}x),\qquad
 h(x)=\begin{cases}\overline{A(x)}/|A(x)|,&A(x)\ne0,\\0,&A(x)=0.\end{cases}
\]
Then $|h|\leq1$ and $hA=|A|$. Let $L=N^d$ and $R_i(y)=P_i(y-1)$.
Measure preservation gives
\[
 \int_X hA\,d\nu
 =\int_X\frac1L\sum_{a=1}^L h(T^ax)
          \frac1N\sum_{n=1}^N g(n)
          \prod_i f_i(T^{a+R_i(n)}x)\,d\nu(x).
\]
For fixed $x$, use $F_0(a)=h(T^ax)\1_{[L]}(a)$ and restrict each
$F_i(b)=f_i(T^bx)$ to a fixed multiple of $[-N^d,N^d]$, extending it by zero.
A coefficient-sum bound chooses that multiple from $\mathbf P$ alone and
contains every $a+R_i(n)$ used in the sum. The finite expression has exactly
the normalization $L^{-1}N^{-1}=N^{-d-1}$. Apply the proved homogeneity form
of TV and integrate the uniform bound on the probability space. For the
original essentially bounded functions, the same pointwise estimate holds
on the conull set where every integer iterate is bounded.
\end{proof}
\ifDetailed
\subsection*{The precise transference obligations}
The checked declaration is \texttt{orbitVNInput\_\allowbreak{}of\_\allowbreak{}finiteVNInput}.
This is the part of the related sequences paper that is needed by the prime
paper. The other paper's main recurrence and pointwise theorems are not inputs.
The preserved source is \texttt{manuscript/}\allowbreak\texttt{related-original/}\allowbreak\texttt{PSSequences.tex};
its label \texttt{p:orbit} is Proposition 1.4. The Lean target is the real-weight,
essentially $1$-bounded specialization required here, rather than a claim to
formalize all variants of that proposition or the entire related paper.

Translation $y\mapsto y-1$ preserves the natural degree of every nonconstant
integer polynomial, including those with nonzero constant term. It therefore
preserves strict degree increase, positive degree and the same maximum degree.
For $1\leq n\leq N$ and $1\leq a\leq N^d$, the elementary coefficient bound
\[
 |a+R_i(n)|\leq C_0(\mathbf P)N^d
\]
holds with $C_0\geq1$ independent of $N$. This is a genuine bound for the
integer support scale used by TV, not an asymptotic support assertion. The
spatial cutoff for $F_0$ is an exact finite interval. Its actual support can
be smaller when the multiplier vanishes; omitted terms are zero, so the
unrestricted integer sum still equals the complete interval average.

The implementation keeps the original representatives throughout. The
weighted orbit average is a finite sum of a.e. strongly measurable functions.
A single conull set controls all integer iterates of all the inputs, so at
each point of that set the entire spatial sample is $1$-bounded. This gives
the finite-form estimate almost everywhere. The same essential bound proves
integrability on a probability space; no representative-change assumption is
needed in the transference theorem. The
explicit multiplier is measurable in the required a.e. sense, bounded by one,
and is set to zero on the zero set of $A$. Thus no measurable-selection theorem
or unproved duality witness is required. The products in the invariant integral
are integrable, and finite summation commutes with integration. The action law
identifies $T^{R_i(n)}T^a$ with $T^{a+R_i(n)}$, also for negative shifts.

The same exponent $K$ selected from $d$ by TV is used throughout. The support
constant is selected from the fixed polynomial family, and TV's bound constant
is then selected before the probability space, action, inputs, weight,
$N,B,E$. No assumption of a discrepancy-specific power saving, an ergodic
limit, or a uniform theorem over changing coefficient families is introduced.
The every-phase formulation is already the interface of TV; its published
supremum translation TV-B remains external, while no separate orbit supremum
bridge is needed for this deduction.
\fi
